% !TEX root = ../../../main.tex
% 来源：中学数学实验教材·第二册（几何）
% 章节：直线形 / 相似形
% 原始文件：3.tex，行 5666-5689
% 类型：tikz
% ---
    \begin{tikzpicture}
\begin{scope}
    \tkzDefPoints{0/0/A,6/0/B, 4/4/C}
    \tkzDrawPolygon(A,B,C)

\tkzDefPointWith[linear, K=.333](A,B)  \tkzGetPoint{B'}
\tkzDefPointWith[linear, K=.333](A,C)  \tkzGetPoint{C'}
\tkzDefPointWith[linear, K=.333](B,C)  \tkzGetPoint{A'}
\tkzDefPointWith[linear, K=.667](A,B)  \tkzGetPoint{B''}
\tkzDefPointWith[linear, K=.667](A,C)  \tkzGetPoint{C''}
\tkzDefPointWith[linear, K=.667](B,C)  \tkzGetPoint{A''}
\tkzDrawSegments[dashed](C',A' B',A'' C'',B'' B'',A' C'',A'')
\tkzDrawSegments(C',B')
    \tkzLabelPoints[below](B',B)
    \tkzLabelPoints[above](C',C)
\tkzLabelPoint[below](A){$A(A')$}
\end{scope}
\begin{scope}[xshift=7cm, scale=.333]
    \tkzDefPoints{0/0/A',6/0/B', 4/4/C'}
    \tkzDrawPolygon(A',B',C')
    \tkzLabelPoints[below](A',B')
    \tkzLabelPoints[above](C')
\end{scope}        
    \end{tikzpicture}
